%% kapa: exam=1 round=32 year=2021 period=2 form=A subject=appraisal_law subject_ko=감정평가관계법규 seq=22 number=22
%% source: Q-Net 감정평가사 기출문제 (cst00309) artlSeq=5208397 file=extracted/1008/5208397/A형/2021년도 제32회 감정평가사 1차 2교시 A형.pdf
%% conversion: kordoc 4.24.51 -> tools/gen1.py (auto); review=proofread-sonnet; flags=none
\begin{question}{22}
국유재산법령상 행정재산의 사용허가에 관한 설명으로 옳은 것은?
\begin{choices}
  \item 중앙관서의 장은 보존용 행정재산의 용도나 목적에 장애가 되지 아니하는 범위에서만 그에 대한 사용허가를 할 수 있다.
  \item 행정재산을 주거용으로 사용허가를 하는 경우에는 일반경쟁의 방법으로 사용허가를 받을 자를 결정하여야 한다.
  \item 중앙관서의 장은 사용허가한 행정재산을 지방자치단체가 직접 공공용으로 사용하기 위하여 필요하게 된 경우에는 그 허가를 철회할 수 있다.
  \item 행정재산으로 할 목적으로 기부를 받은 재산에 대하여 기부자에게 사용허가하는 경우에는 그 사용허가기간은 5년 이내로 한다.
  \item 행정재산의 사용허가에 관하여는 「국유재산법」에서 정한 것을 제외하고는「민법」의 규정을 준용한다.
\end{choices}
\end{question}
\answer{3}
